\part{Picard 概形：存在定理}\label{fga3.v}\thispagestyle{fancy}\begin{original}{}\label{fga3.v-original-citation}\skipsinglepar\par{}A. Grothendieck．“Technique de descente et théorèmes d'existence en géométrie algébrique, V: Les schémas de Picard: Théorèmes d'existence”．\emph{Séminaire Bourbaki} \textbf{14}（1961–62），第 232 次报告，1962 年 2 月。\url{http://www.numdam.org/book-part/SB_1961-1962__7__143_0/}\end{original}\setcounter{chapter}{0}\chapter{相对 Picard 群与函子}\thispagestyle{fancy}\label{fga3.v-1}\par{}对于每个预概形（更一般地，每个赋环空间）\(X\)，定义 \(X\) 的（\emph{绝对}）\emph{Picard 群}，记为 \(\operatorname {Pic}(X)\)，为 \(X\) 上可逆模（即局部同构于 \(\mathcal {O}_X\) 的模）的同构类所成的群。因此有典范同构 \begin{eqn}{1.1}\setnumber{1.1}\label{fga3.v-1-equation-1.1}\skipsinglepar\[       \operatorname {Pic}(X) \xrightarrow {\sim } \operatorname {H}^1(X,\mathcal {O}_X^\times )     \tag{1.1}     \]\end{eqn}，其中 \(\mathcal {O}_X^\times \) 表示 \(\mathcal {O}_X\) 的可逆元层（可将它等同于可逆模 \(\mathcal {O}_X\) 的自同构层）。注意，\(X\mapsto \operatorname {Pic}(X)\) 显然是关于 \(X\) 的反变函子，且同构 \hyperref[fga3.v-1-equation-1.1]{式~\ref*{fga3.v-1-equation-1.1}} 是函子性的。\par{}设 \(X\) 是预概形 \(S\) 上的预概形。由以上讨论，当 \(S'\) 在 \(S\) 上预概形的范畴 \(\mathtt {Sch}_{/S}\) 中变化时，就得到反变函子 \(S'\mapsto \operatorname {Pic}(X\times _S S')\)。这个函子不可能“可表”（\hyperref[fga3.ii-a.1]{FGA 3.II, §A}），因为我们拟加以分类的可逆模具有 \emph{非平凡自同构}，因而此函子并不具有“\emph{局部性质}”（[\hyperref[Gro1960a]{Gro1960a}，IV，5.4]）。因此，有必要使其“局部化”：对每个相对预概形 \(X/S\)，引入一个相对性质的群 \begin{eqn}{1.2}\setnumber{1.2}\label{fga3.v-1-equation-1.2}\skipsinglepar\[       \operatorname {Pic}'(X/S)       = \operatorname {H}^0(S,\operatorname {R}^1f_*(\mathcal {O}_X^\times ))     \tag{1.2}     \]\end{eqn}（其中 \(f\colon  X\to  S\) 是结构态射；参见 \hyperref[fga3.ii-c.3]{FGA 3.II, §C.3}）。在 \emph{同上} 中，这个群称为相对 Picard 群；不过，出于下文将说明的原因，这里最好称之为 \(X/S\) 的 \emph{受限相对 Picard 群}。当 \(S'\) 在 \(\mathtt {Sch}_{/S}\) 中变化时，\(S'\mapsto \operatorname {Pic}'(X\times _S S'/S')\) 是关于 \(S'\) 的反变函子，也记为 \(\mathscr {P}\kern  -1pt\mathit {ic}'_{X/S}\)，其本质上由公式 \begin{eqn}{1.3}\setnumber{1.3}\label{fga3.v-1-equation-1.3}\skipsinglepar\[       \mathscr {P}\kern  -1pt\mathit {ic}'_{X/S}(S')       = \operatorname {Pic}(X\times _S S'/S').     \tag{1.3}     \]\end{eqn} 给出。现在，这个函子具有“局部性质”，因为我们已经作了为此所需的处理。直观地说，\hyperref[fga3.v-1-equation-1.3]{式~\ref*{fga3.v-1-equation-1.3}} 的右端可理解为以参数预概形 \(S'/S\) 为指标的、\(X/S\)（的纤维）上可逆层的类的“代数族”所成的集合。当函子 \(\mathscr {P}\kern  -1pt\mathit {ic}'\) 可表时，将表示它的 \(S\) 上的预概形记为 \(\underline {\operatorname {Pic}}_{X/S}\)，称为 \(X\) 在 \(S\) 上的 \emph{Picard 预概形}；这时有 \begin{eqn}{1.4}\setnumber{1.4}\label{fga3.v-1-equation-1.4}\skipsinglepar\[       \operatorname {Hom}_S(S',\underline {\operatorname {Pic}}_{X/S})       \cong  \mathscr {P}\kern  -1pt\mathit {ic}'_{X/S}(S')       = \operatorname {Pic}'(X\times _S S'/S').     \tag{1.4}     \]\end{eqn}。然而，在某些重要情形中，\(\mathscr {P}\kern  -1pt\mathit {ic}'_{X/S}\) 不可表（例如域 \(k\) 上没有 \(k\)-有理点的“Brauer–Severi”簇），但相对 Picard 预概形仍有自然的定义。这是因为，从绝对 Picard 群 \(\operatorname {Pic}(X\times _S S'/S')\) 定义函子 \(\mathscr {P}\kern  -1pt\mathit {ic}'\) 时，我们所作的局部化还不够；更确切地说，\(\mathscr {P}\kern  -1pt\mathit {ic}'\) 一般并不“与忠实平坦下降相容”。下面详细说明。\par{}令 \((\mathscr {M})\) 为所有 \emph{忠实平坦且拟紧} 的预概形态射所成的集合；此集合在基变换及复合下稳定。设 \(P\) 是从 \(\mathtt {Sch}_{/S}\) 到集合范畴的反变函子。对每个满足 \(u\in (\mathscr {M})\) 的 \(S\)-态射 \(u\colon  T'\to  T\)，考虑图表 \begin{eqn}{1.5}\setnumber{1.5}\label{fga3.v-1-equation-1.5}\skipsinglepar\[       P(T)       \to  P(T')       \rightrightarrows  P(T'\times _T T')     \tag{1.5}     \]\end{eqn}；它由对图表 \[     T     \leftarrow  T'     \underset {\mathrm {pr}_2}{\overset {\mathrm {pr}_1}{\leftleftarrows }} T'\times _T T'.   \] 施以 \(P\) 得到。若 \(P\) 可表，则由下降理论（\hyperref[fga3.i-b.1-theorem-2]{FGA 3.I，§B，定理 2}）可知，对所有 \(u\in (\mathscr {M})\)，图表 \hyperref[fga3.v-1-equation-1.5]{式~\ref*{fga3.v-1-equation-1.5}} 都正合。我们称这一事实为 \(P\) 与 \((\mathscr {M})\) 相容；就此处而言，也就是说 \(P\)“与忠实平坦下降相容”，或说 \(\mathtt {Sch}_{/S}\) 上的 \emph{“预层”} \(P\) 对集合 \((\mathscr {M})\) 所给出的局部化概念而言是一个 \emph{“层”}。对于任意 \(P\)，用一个在通常拓扑局部化情形中熟知的标准过程，可以给它配备一个“层”\(\mathscr {P}\) 及函子同态 \(P\to \mathscr {P}\)，后者具有显然的泛性质。\(\mathscr {P}\) 的构造可明确写成如下形式：为定义 \(\mathscr {P}(T)\)，对每个满足态射 \(u\colon  T'\to  T\) 属于 \((\mathscr {M})\) 的、\(T\) 上的 \(T'\)，用 \(\overline {\operatorname {H}}^0(T'/T,P)\) 表示 \(P(T')\) 的如下子集：其元素 \(\xi \) 在 \(P(T'\times _T T')\) 中的像 \(\xi _1,\xi _2\) 满足，存在 \((\mathscr {M})\) 中的态射 \(v\colon  T''\to  T'\times _T T'\)，使 \(\xi _1\) 与 \(\xi _2\) 在 \(P(T'')\) 中有相同的像。\par{}（注意：如此定义的集合 \(\overline {\operatorname {H}}^0\) 比 \hyperref[fga3.i-a.4.a]{FGA 3.I, §A.4.a} 中引入的集合 \(\operatorname {H}^0(T'/T,P)\) 大。）让 \(T'\) 在固定的 \(T\) 上变化（始终要求 \(u\in (\mathscr {M})\)），各个 \(\overline {\operatorname {H}}^0(T'/T,P)\) 构成一个归纳系（在所有 \(T'\) 构成的集合上，以支配关系给出预序），并令 \begin{eqn}{1.6}\setnumber{1.6}\label{fga3.v-1-equation-1.6}\skipsinglepar\[       \mathscr {P}(T)       = \varinjlim _{T'} \overline {\operatorname {H}}^0(T'/T,P).     \tag{1.6}     \]\end{eqn}。这个表达式关于 \(T\) 的函子性是显然的。\par{}当 \[     P(T)     = \operatorname {Pic}(X\times _S T)   \] 时，称由 \hyperref[fga3.v-1-equation-1.6]{式~\ref*{fga3.v-1-equation-1.6}} 定义的 \(\mathtt {Sch}_{/S}\) 上的反变函子为 \(X\) 在 \(S\) 上的 \emph{相对 Picard 函子}，记为 \(\mathscr {P}\kern  -1pt\mathit {ic}_{X/S}\)；并将群 \(\mathscr {P}\kern  -1pt\mathit {ic}_{X/S}(S)\) 定义为 \(X\) 在 \(S\) 上的 \emph{相对 Picard 群}，记为 \(\operatorname {Pic}(X/S)\)。于是有显然的双射 \begin{eqn}{1.7}\setnumber{1.7}\label{fga3.v-1-equation-1.7}\skipsinglepar\[       \mathscr {P}\kern  -1pt\mathit {ic}_{X/S}(T)       \xrightarrow {\sim } \operatorname {Pic}(X\times _S T/T).     \tag{1.7}     \]\end{eqn}。因此，\(\operatorname {Pic}(X/S)\) 的一个元素由群 \(\operatorname {Pic}(X\times _S S')\) 的一个元素 \(\xi '\) 给出（其中 \(S'\to  S\) 忠实平坦且拟紧），并要求能够找到忠实平坦拟紧态射 \(S''\to  S'\times _S S'\)，使 \(\xi '\) 在 \(\operatorname {Pic}(X\times _S S'')\) 中的两个逆像相同。满足刚才条件的元素 \(\xi '\)（属于 \(\operatorname {Pic}(X\times _S S')\)）和 \(\xi _1\)（属于 \(\operatorname {Pic}(X\times _S S_1)\)）定义 \(\operatorname {Pic}(X/S)\) 的同一元素，当且仅当存在忠实平坦拟紧态射 \(S'_1\to  S'\times _S S_1\)，使这两个元素在 \(\operatorname {Pic}(X\times _S S'_1)\) 中的像相等。实际运算时，往往改用上面引入的函子 \(P'=\mathscr {P}\kern  -1pt\mathit {ic}'_{X/S}\) 更为方便。立刻可见，典范态射 \(P\to  P'\) 定义了一个 \emph{同构} \begin{eqn}{1.8}\setnumber{1.8}\label{fga3.v-1-equation-1.8}\skipsinglepar\[       \mathscr {P}       \xrightarrow {\sim } \mathscr {P}'     \tag{1.8}     \]\end{eqn}，从而用 \(\mathscr {P}\kern  -1pt\mathit {ic}'_{X/S}=P'\) 对 \(\mathscr {P}\kern  -1pt\mathit {ic}_{X/S}\) 给出了一种通常更方便的描述。由下面的 \hyperref[fga3.v-2-corollary-2.3]{推论~\ref*{fga3.v-2-corollary-2.3}}，若将刚才对 \(\operatorname {Pic}(X/S)\) 的描述中的 \(P\) 换为 \(P'\)，则至少在 \emph{同上} 所给的条件下，可以取 \(S''=S'\times _S S'\) 及 \(S'_1=S'\times _S S_1\)。\par{}若函子 \(\mathscr {P}\kern  -1pt\mathit {ic}_{X/S}\) 可表，就说 \(X/S\) 有 Picard 预概形，并称表示此函子的 \(S\) 上的预概形为 \(X\) 在 \(S\) 上的 \emph{Picard 预概形}，记为 \(\underline {\operatorname {Pic}}_{X/S}\)。显然，为此只须 \(P'=\mathscr {P}\kern  -1pt\mathit {ic}_{X/S}\) 可表，因为此时 \(P'\) 已是一个“层”，而等式 \hyperref[fga3.v-1-equation-1.8]{式~\ref*{fga3.v-1-equation-1.8}} 表明，态射 \(P'\to \mathscr {P}'\) 可等同于典范态射 \begin{eqn}{1.9}\setnumber{1.9}\label{fga3.v-1-equation-1.9}\skipsinglepar\[       \mathscr {P}\kern  -1pt\mathit {ic}'_{X/S}       \to  \mathscr {P}\kern  -1pt\mathit {ic}_{X/S}     \tag{1.9}     \]\end{eqn}，后者于是是同构。这说明，我们的术语与上面在 \hyperref[fga3.v-1-equation-1.4]{式~\ref*{fga3.v-1-equation-1.4}} 中引入的术语相容。一般地，只要 \(\underline {\operatorname {Pic}}_{X/S}\) 存在，它就由函子性同构 \begin{eqn}{1.10}\setnumber{1.10}\label{fga3.v-1-equation-1.10}\skipsinglepar\[       \operatorname {Hom}_S(S',\underline {\operatorname {Pic}}_{X/S})       \xrightarrow {\sim } \operatorname {Pic}(X\times _S S'/S').     \tag{1.10}     \]\end{eqn} 所定义。\chapter{各类相对与绝对 Picard 群之间的关系}\thispagestyle{fancy}\label{fga3.v-2}\begin{proposition}{2.1}\setnumber{2.1}\label{fga3.v-2-proposition-2.1}\skipsinglepar\par{}设 \(f\colon  X\to  S\) 是满足 \(\mathcal {O}_S\xrightarrow {\sim } f_*(\mathcal {O}_X)\) 的态射。则有正合列 \[       0       \to  \operatorname {Pic}(S)       \to  \operatorname {Pic}(X)       \to  \operatorname {Pic}'(X/S).     \]。若 \(X\) 在 \(S\) 上有截面，则最后一个态射满射，即有同构 \[       \operatorname {Pic}'(X/S)       \xrightarrow {\sim } \operatorname {Pic}(X)/\operatorname {Pic}(S).     \]。\begin{proof}\par{}此正合列可视为对应于 \(f\) 和 \(\mathcal {O}_X\) 的 Leray 谱序列所给出的正合列的低次部分。第二个断言同样是形式的。\end{proof}\end{proposition}\begin{proposition}{2.2}\setnumber{2.2}\label{fga3.v-2-proposition-2.2}\skipsinglepar\par{}设 \(f\colon  X\to  S\) 是满足 \(\mathcal {O}_S\xrightarrow {\sim } f_*(\mathcal {O}_X)\) 的拟紧分离态射，\(S'\to  S\) 是忠实平坦拟紧态射。那么：\begin{enumerate}
      
 \item[i.]\(\operatorname {Pic}'(X/S)\to \operatorname {Pic}'(X\times _S S'/S')\) 是单射；

      
 \item[ii.]若 \(X\) 在 \(S\) 上局部有截面（即每个 \(s\in  S\) 都有开邻域 \(U\)，使得 \(X|U\) 在 \(U\) 上有截面），则图表 \[           \operatorname {Pic}'(X/S)           \to  \operatorname {Pic}'(X\times _S S'/S')           \rightrightarrows  \operatorname {Pic}'(X\times _S S''/S'')         \]（其中 \(S''=S'\times _S S'\)）正合。

    \end{enumerate}\begin{proof}\par{}借助忠实平坦下降的初等性质，第一个断言由下面的一般注记得到。若 \(f\colon  X\to  S\) 是满足 \(\mathcal {O}_S\xrightarrow {\sim } f_*(\mathcal {O}_X)\) 的态射，则从 \(S\) 上有限型局部自由模的范畴到 \(X\) 上有限型局部自由模的范畴的函子 \(\mathscr {F}\mapsto  f^*(\mathscr {F})\) 是全忠实的；其本质像是 \(X\) 上满足如下条件的模 \(\mathscr {G}\)：\(f_*(\mathscr {G})\) 局部自由，且典范同态 \[         f^*f_*(\mathscr {G})         \to  \mathscr {G}       \] 为同构。第二个断言已在 \hyperref[fga3.i-b.4]{FGA 3.I, §B.4} 中借助下降理论证明。\end{proof}\end{proposition}\par{}\hyperref[fga3.v-2-proposition-2.2]{命题~\ref*{fga3.v-2-proposition-2.2}} 的结果也可表述如下：\begin{corollary}{2.3}\setnumber{2.3}\label{fga3.v-2-corollary-2.3}\skipsinglepar\par{}在 \hyperref[fga3.v-2-proposition-2.2]{命题~\ref*{fga3.v-2-proposition-2.2}} 的条件下，典范同态 \hyperref[fga3.v-1-equation-1.9]{式~\ref*{fga3.v-1-equation-1.9}} \(\mathscr {P}\kern  -1pt\mathit {ic}'_{X/S}\to \mathscr {P}\kern  -1pt\mathit {ic}_{X/S}\) 是单射；若 \(X\) 局部地在 \(S\) 上有截面，则它还是双射。（后一情形中，相对 Picard 群 \(\operatorname {Pic}(X/S)\) 等同于受限相对 Picard 群 \(\operatorname {Pic}'(X/S)\)。）\end{corollary}\par{}将它与 \hyperref[fga3.v-2-proposition-2.1]{命题~\ref*{fga3.v-2-proposition-2.1}} 结合，便得到：\begin{corollary}{2.4}\setnumber{2.4}\label{fga3.v-2-corollary-2.4}\skipsinglepar\par{}在 \hyperref[fga3.v-2-proposition-2.2]{命题~\ref*{fga3.v-2-proposition-2.2}} 的条件下，有正合列 \[       0       \to  \operatorname {Pic}(S)       \to  \operatorname {Pic}(X)       \to  \operatorname {Pic}(X/S).     \]。若 \(X\) 在 \(S\) 上有截面，则最后一个同态满射，即有同构 \[       \operatorname {Pic}(X/S)       \xrightarrow {\sim } \operatorname {Pic}(X)/\operatorname {Pic}(S).     \]。\end{corollary}\begin{remark}{2.5}\setnumber{2.5}\label{fga3.v-2-remark-2.5}\skipsinglepar\par{}设 \(f\colon  X\to  S\) 是满足 \(\mathcal {O}_S\xrightarrow {\sim } f_*(\mathcal {O}_X)\) 的态射，\(g\) 是 \(X\) 在 \(S\) 上的一个截面。设 \(\mathscr {L}\) 是 \(X\) 上的可逆模。将 \(\mathscr {L}\) 的 \emph{\(g\)-刚性化} 定义为一个同构 \(\mathcal {O}_S\xrightarrow {\sim } g^*(\mathscr {L})\)，并将 \emph{\(g\)-刚性化可逆模} 定义为 \(X\) 上配有一个 \(g\)-刚性化的可逆模 \(\mathscr {L}\)。这样一个结构的每个自同构都是平凡的，而 \(\operatorname {Pic}'(X/S)\) 可等同于 \(S\) 上 \(g\)-刚性化可逆模的同构类所成的集合。（正是这一事实使我们能利用下降理论证明 \hyperref[fga3.v-2-proposition-2.2]{命题~\ref*{fga3.v-2-proposition-2.2}} 的 (ii)。）这给出了对 \(\operatorname {Pic}(X/S)\) 的新解释，至少当 \(f\) 还拟紧且分离时是如此；因为这时由 \hyperref[fga3.v-2-corollary-2.3]{推论~\ref*{fga3.v-2-corollary-2.3}} 可知 \(\operatorname {Pic}(X/S)\xrightarrow {\sim }\operatorname {Pic}'(X/S)\)。\end{remark}\begin{remark}{2.6}\setnumber{2.6}\label{fga3.v-2-remark-2.6}\skipsinglepar\par{}设 \(f\colon  X\to  S\) 是 \hyperref[fga3.v-2-proposition-2.2]{命题~\ref*{fga3.v-2-proposition-2.2}} 中的态射，\(S'\to  S\) 是忠实平坦拟紧态射，使得存在一个 \(S\)-态射 \(S'\to  X\)，即 \(X'=X\times _S S'\) 在 \(S'\) 上有截面。令 \(S''=S'\times _S S'\)、\(X''=X\times _S S''\)，并考虑正合列 \[       \operatorname {Pic}(X/S)       \to  \operatorname {Pic}(X'/S')       \rightrightarrows  \operatorname {Pic}(X''/S'').     \]。施以 \hyperref[fga3.v-2-corollary-2.4]{推论~\ref*{fga3.v-2-corollary-2.4}}，得到正合列 \[       \operatorname {Pic}(X/S)       \to  \operatorname {Pic}(X')/\operatorname {Pic}(S')       \rightrightarrows  \operatorname {Pic}(X'')/\operatorname {Pic}(S'').     \]。特别地，相对 Picard 群的每个元素都已经“来自”\(\operatorname {Pic}(X')\) 的某个元素。这大大简化了上一节对相对 Picard 群的描述，甚至也简化了 \(X\) 在 \(S\) 上的 Picard 函子的描述，因为对每个 \(S\) 上的 \(T\)，都可以将上述讨论用于 \(X\times _S T/T\) 及态射 \(T'=S'\times _S T\to  T\)。例如，若 \(f\) 本身忠实平坦，就可取 \(S'=X\)；于是，只要 \(f\) 还为有限型（分别为光滑态射，等等），在描述相对 Picard 函子 \(T\mapsto \operatorname {Pic}(X\times _S T/T)\) 时，就可以只考虑有限型（分别为光滑态射，等等）的基变换 \(T'\to  T\)。若 \(f\) 射影且平坦，\(S\) 局部 Noether，则可以证明，上面可取这样的 \(S'\to  S\)，使 \(S'\) 是 \(S\) 中覆盖 \(S\) 的若干开集 \(S_i\) 的 \emph{平坦覆盖} \(S'_i\) 的直和；若 \(f\) 还可分，则可取 \(S'_i\) 在 \(S_i\) 上 étale。\end{remark}\chapter{主要存在定理：陈述}\thispagestyle{fancy}\label{fga3.v-3}\par{}关于 Picard 预概形的存在性，目前还没有一个涵盖所有已知情形的陈述，甚至连猜想也没有。如果可以这样说，一个“实际上必要”的条件是 \(f\colon  X\to  S\) 为 \emph{固有}（以保证基本的有限性性质）并且 \emph{平坦}。这些条件并不充分，即使 \(S\) 是域 \(k\)（譬如复数域 \(\mathbb {C}\)）上的对偶数代数 \(k[t]/(t^2)\) 的谱，且 \(X\) 的维数为 \(1\)，也是如此。在撰写本报告时，关于 Picard 预概形的最重要的存在定理都由下列定理推出：\begin{theorem}{3.1}\setnumber{3.1}\label{fga3.v-3-theorem-3.1}\skipsinglepar\par{}设 \(f\colon  X\to  S\) 是局部 Noether 预概形之间的态射。假设：\begin{enumerate}
      
 \item[i.]\(f\) 射影；

      
 \item[i.]\(f\) 平坦；

      
 \item[i.]\(f\) 的几何纤维整。

    \end{enumerate}在这些条件下，\(\underline {\operatorname {Pic}}_{X/S}\) 存在。\end{theorem}\par{}下面两节将给出的证明梗概也将同时说明下述事实。设 \(\xi \) 是 \(\underline {\operatorname {Pic}}_{X/S}\) 的一个截面，对应于 \(X/S\) 上的一个甚丰沛层 \(\mathcal {O}_X(1)\)（即由射影嵌入 \(X\to \mathbb {P}(\mathscr {E})\) 诱导的层）；那么，\(\underline {\operatorname {Pic}}_{X/S}\) 中存在开子集 \(U\)，它是 \(S\) 上拟射影开子集的不交并，使得 \(U\) 在由 \(\xi \) 给出的平移下稳定，且 \(\underline {\operatorname {Pic}}_{X/S}\) 是开集 \(U\setminus  n\xi \) 的递增并（每个开集都同构于 \(U\)）。特别地，由此可得，在 \hyperref[fga3.v-3-theorem-3.1]{定理~\ref*{fga3.v-3-theorem-3.1}} 的条件下，\(\underline {\operatorname {Pic}}_{X/S}\) 在 \(S\) 上 \emph{分离}。\begin{remark}{3.2}\setnumber{3.2}\label{fga3.v-3-remark-3.2}\skipsinglepar\par{}一些例子（例如取 \(S\) 为离散赋值环的谱，\(X\) 在 \(S\) 上的相对维数为 \(1\)）表明，若删去 \hyperref[fga3.v-3-theorem-3.1]{定理~\ref*{fga3.v-3-theorem-3.1}} 中的假设 (iii)，代之以较弱的假设：对每个 \(s\in  S\)，同态 \(k(s)\to \operatorname {H}^0(X_s,\mathcal {O}_{X_s})\) 均为同构，那么 \(\underline {\operatorname {Pic}}_{X/S}\) 在 \(S\) 上未必分离。这既可能出现在 \(f\) 的几何纤维约化、但一个整的泛几何纤维经特殊化“分裂”为两个不可约分支时，也可能出现在 \(f\) 的几何纤维不可约、但一个整的泛几何纤维特殊化为一个“多重纤维”时。第一种情形的例子是一条圆锥曲线退化为两条相交直线；D. Mumford 曾给我看过第二种情形的一个例子：一条椭圆曲线退化为一条二重椭圆曲线。这些例子在任意特征下都成立。\end{remark}\begin{remark}{3.3}\setnumber{3.3}\label{fga3.v-3-remark-3.3}\skipsinglepar\par{}在 \hyperref[fga3.v-3-theorem-3.1]{定理~\ref*{fga3.v-3-theorem-3.1}} 的条件下，我不知道 \(\underline {\operatorname {Pic}}_{X/S}\) 是否是若干开集的不交并，其中每个开集在 \(S\) 上为有限型，因而拟射影。注意，与 Hilbert 概形的情形（\hyperref[fga3.iv]{FGA~\ref*{fga3.iv}}）相同，对 Hilbert 多项式 \(Q\in \mathbb {Q}[t]\) 的研究可以将 \(\underline {\operatorname {Pic}}_{X/S}\) 分解为开集 \(\underline {\operatorname {Pic}}_{X/S}^Q\) 的不交和；这些开集在 \(S\) 上似乎应当拟射影。至少当 \(f\) 为光滑态射时，下一次报告将证明这一点。请注意，若将假设 (i) 换为“\(X\) 在 \(S\) 上 \emph{局部地} 射影”（这足以证明 \hyperref[fga3.v-3-theorem-3.1]{定理~\ref*{fga3.v-3-theorem-3.1}}，因为 \(\underline {\operatorname {Pic}}_{X/S}\) 的存在问题显然在 \(S\) 上是局部的），则很容易给出满足 \emph{\(\underline {\operatorname {Pic}}_{X/S}\) 含有在 \(S\) 上并非有限型的连通分支} 的例子。例如，设 \(X_0\) 是代数闭域 \(k\) 上的非奇异射影代数簇，配有自同构 \(u\) 及 \(X_0\) 的 Néron–Severi 群中的元素 \(\xi \)，使得各 \(u^n(\xi )\) 两两不同。譬如，可取 \(X_0\) 为一条椭圆曲线 \(E\) 与自身的乘积，取 \(u\) 为 \(E\times  E\) 的自同构 \((x,y)\mapsto (x,y+x)\)。设 \(S\) 是两条非奇异不可约曲线的并，它们交于两点 \(a\) 和 \(b\)。在 \(S\) 上存在群 \(\mathbb {Z}\) 的连通主覆盖 \(P\)；利用 \(u\) 所定义的 \(\mathbb {Z}\) 在 \(X_0\) 上的作用，便得到 \(S\) 上纤维为 \(X_0\) 的伴随纤维丛（在 \(S\setminus \{a\}\) 和 \(S\setminus \{b\}\) 上平凡）；在这里所考虑的特殊情形中，它实际上是 \(S\) 上的一个 \emph{Abel} 概形。容易看出，\(\underline {\operatorname {Pic}}_{X/S}\) 也是由 \(P\) 及 \(\mathbb {Z}\) 通过 \(u\) 在 \(\underline {\operatorname {Pic}}_{X_0/k}\) 上的作用所得到的伴随纤维丛；它含有一个同构于 \(P\times \underline {\operatorname {Pic}}_{X_0/k}^0\) 的连通分支（其中 \(\underline {\operatorname {Pic}}^0\) 表示 \(\underline {\operatorname {Pic}}\) 中单位元的连通分支），这个分支在 \(S\) 上并非有限型。（在 \hyperref[fga3.v-3-remark-3.3]{注~\ref*{fga3.v-3-remark-3.3}} 所述的 \(S\) 上非分离 Picard 预概形的各种情形中，也可以构造类似现象。）\emph{[编者]} 这里提出的问题已由 Mumford 给出肯定回答（见 \hyperref[fga3.vi]{FGA~\ref*{fga3.vi}}）。\end{remark}\chapter{相对 Cartier 除子与射影丛}\thispagestyle{fancy}\label{fga3.v-4}\par{}我们只需用到正除子，因此本节以下省略“正”这一限定词。\par{}设 \(X\) 是一个预概形。\(X\) 上的 \emph{Cartier 除子}，简称除子，是 \(X\) 的一个闭子预概形 \(D\)，其定义理想 \(\mathscr {J}\) 是一个 \emph{可逆} 模，即局部由 \(\mathcal {O}_X\) 的一个 \emph{非零因子} 截面生成。给 \(D\) 配上可逆模 \[     \mathscr {L}(D)     = \mathscr {J}^{-1}   \]，则典范单射 \(\mathscr {J}\to \mathcal {O}_X\) 给出典范同态 \[     s_D\colon  \mathcal {O}_X     \to  \mathscr {J}^{-1}     = \mathscr {L}(D)   \]，即 \(s_D\in \Gamma (X,\mathscr {L}(D))\)。给定一个除子，本质上 \emph{等价} 于给定 \(X\) 上的一个可逆模 \(\mathscr {L}\) 及其处处不是零因子的截面 \(s\)：给这样一对 \((\mathscr {L},s)\) 配上 \(s\) 的“除子”，记为 \(\operatorname {div}(s)\)。对 \(X\) 上给定的可逆模 \(\mathscr {L}\)，定义 \(\mathscr {L}\) 的除子 \(D\) 所成的集合，与商集 \(\Gamma (X,\mathscr {L})^{\times }/\Gamma (X,\mathcal {O}_X^\times )\) 一一对应；其中 \(\Gamma (X,\mathscr {L})^\times \) 表示 \(\Gamma (X,\mathscr {L})\) 中处处不是零因子的截面所成的子集。\par{}现在设有一个局部有限型态射 \(f\colon  X\to  S\)；为简单起见，假设 \(S\) 局部 Noether。设 \(\mathscr {J}\) 是 \(X\) 上的凝聚理想，\(D\) 是它在 \(X\) 中定义的子概形，并令 \(x\in  X\)、\(s=f(x)\)。我们将证明下列条件等价：\begin{enumerate}
    
 \item[i.]\(\mathscr {J}\) 在 \(x\) 处可逆（即 \(\mathscr {J}_x\) 由 \(\mathcal {O}_{X,x}\) 的一个正则元生成），且 \(D\) 在 \(x\) 处在 \(S\) 上平坦。

    
 \item[ii.]\(X\) 和 \(D\) 在 \(x\) 处在 \(S\) 上平坦，且 \(D_s\) 在点 \(x\) 处是纤维 \(X_s\) 上的 Cartier 除子。

    
 \item[iii.]\(X\) 在 \(x\) 处在 \(S\) 上平坦，且 \(\mathscr {J}_x\) 由一个元素 \(f_x\) 生成，该元素在 \(X_s\) 上诱导一个非零因子芽。

  \end{enumerate}此时称 \(D\) 在所考虑的点处为 \(X/S\) 上的一个 \emph{相对 Cartier 除子}，或简称相对除子。注意，在 (i) 中，\(D\) 在 \(x\) 的某个邻域内的各点处也都是相对除子。因此，若 \(X\) 和 \(D\) 在 \(S\) 上平坦，且 \(D\) 在 \(S\) 上固有，那么，使 \(D_s\) 成为 \(X_s\) 中的 Cartier 除子的 \(s\in  S\)（即使 \(D\) 在 \(X_s\) 的各点处都是相对 Cartier 除子的那些点）所成的集合是 \(S\) 的开子集。上面的定义也已作了必要安排，以保证相对 Cartier 除子的概念在任意基变换 \(S'\to  S\) 下稳定。于是考虑 \(X/S\) 上相对除子的集合 \(\operatorname {Div}(X/S)\)，以及由 \[     \mathscr {D}\kern  -1pt\mathit {iv}_{X/S}(S')     = \operatorname {Div}(X\times _S S'/S').   \] 定义的、关于在 \(S\) 中变化的 \(S'\) 的反变函子。假设 \(X\) 在 \(S\) 上平坦且固有。由相对 Cartier 除子的刻画 (ii)，可将 \(\mathscr {D}\kern  -1pt\mathit {iv}_{X/S}\) 看作 \hyperref[fga3.iv]{FGA~\ref*{fga3.iv}} 中定义的函子 \(\mathscr {H}\kern  -1.75pt\mathit {ilb}_{X/S}\) 的子函子；由上述注记，包含态射 \[     \mathscr {D}\kern  -1pt\mathit {iv}_{X/S}     \to  \mathscr {H}\kern  -1.75pt\mathit {ilb}_{X/S}   \] “可由开浸入表示”（参见 [\hyperref[Gro1960a]{Gro1960a}，IV，3.13]）。利用 \hyperref[fga3.iv]{FGA~\ref*{fga3.iv}} 的主要存在定理，得到：\begin{proposition}{4.1}\setnumber{4.1}\label{fga3.v-4-proposition-4.1}\skipsinglepar\par{}假设 \(f\colon  X\to  S\) 射影且平坦。那么函子 \(\mathscr {D}\kern  -1pt\mathit {iv}_{X/S}\) 可表；更确切地说，它由 \(\underline {\operatorname {Hilb}}_{X/S}\) 的一个开集表示。\end{proposition}\par{}给定 \(X/S\) 上的甚丰沛层 \(\mathcal {O}_X(1)\)，利用 \(\underline {\operatorname {Hilb}}_{X/S}\) 按 Hilbert 多项式 \(Q\in \mathbb {Q}[t]\) 典范分解为开集 \(\underline {\operatorname {Hilb}}_{X/S}^Q\) 之和的分解，得到类似的分解 \[     \underline {\operatorname {Div}}_{X/S}     = \sqcup _{Q\in \mathbb {Q}[t]} \underline {\operatorname {Div}}_{X/S}^Q   \]，它是若干两两不交、在 \(S\) 上 \emph{拟射影} 的开集之和。\par{}利用映射 \(D\mapsto \mathscr {L}(D)\)，得到一个函子性同态 \[     \mathscr {D}\kern  -1pt\mathit {iv}_{X/S}     \to  \mathscr {P}\kern  -1pt\mathit {ic}_{X/S}   \tag{+}   \]；我们拟研究它。在相当一般的条件下，它相对可表（[\hyperref[Gro1960a]{Gro1960a}，IV，3]）。因此，从 \(\mathscr {P}\kern  -1pt\mathit {ic}_{X/S}(S')\) 的一个元素 \(\xi \) 出发，并为简化记号而假设 \(S'=S\)，我们将证明 \(\mathscr {D}\kern  -1pt\mathit {iv}_{X/S}\) 的相应子函子可表。首先考虑 \(\xi \) 由 \(X\) 上的一个可逆模 \(\mathscr {L}\) 定义的情形。假设 \(X\) 在 \(S\) 上固有且平坦，\(X\) 在 \(S\) 上的几何纤维都是整的；这还蕴含（[\hyperref[GD1960]{GD1960}，III，§7]）\(\mathcal {O}_S\xrightarrow {\sim } f_*(\mathcal {O}_X)\)，而且这一等式在任意基变换 \(S'\to  S\) 后仍成立。于是，\(X/S\) 上满足 \(\mathscr {L}(D)\) 与 \(\mathscr {L}\) 定义 \(\operatorname {Pic}(X/S)=\mathscr {P}\kern  -1pt\mathit {ic}_{X/S}(S)\) 中同一元素的相对 Cartier 除子 \(D\)，也就是由 \hyperref[fga3.v-2-corollary-2.4]{推论~\ref*{fga3.v-2-corollary-2.4}} 可知满足 \(\mathscr {L}(D)\) 与 \(\mathscr {L}\) 在 \(S\) 上局部同构的相对 Cartier 除子，与商层 \(f_*(\mathscr {L})^{\times }/\mathscr {O}_S^\times \) 的截面一一对应。此对应与基变换相容。\emph{同上} 及 [Mumford–Tate seminar，1962] 中的“Künneth”型一般论证表明，\(X/S\) 的固有性及 \(\mathscr {L}\) 在 \(S\) 上的平坦性，蕴含 \(S\) 上存在一个凝聚模 \(\mathscr {Q}\)（在唯一同构意义下确定）及层的同构 \[     f_*(\mathscr {L})     \xrightarrow {\sim } \mathscr {H}\kern  -2.5pt\mathit {om}_{\mathcal {O}_X}(\mathscr {Q},\mathcal {O}_S)   \]；此外，\(\mathscr {Q}\) 的构造与基变换相容。这里，\(f_*(\mathscr {L})^\times \) 表示 \(f_*(\mathscr {L})\) 的集合子层，其在 \(U\) 上的截面是 \(\mathscr {L}\) 在 \(f^{-1}(U)\) 上那些在 \(f^{-1}(U)/U\) 上定义相对 Cartier 除子的截面，即在各 \(X_s\) 上（\(s\in  U\)）诱导出非零因子截面的截面。利用纤维 \(X_s\) 整的假设，这一条件等价于在纤维 \(X_s\) 上诱导的截面不恒为零；用局部同态 \(\mathscr {Q}\to \mathcal {O}_S\) 来说，便是这些同态满射（Nakayama）。这说明，\(f_*(\mathscr {L})^{\times }/\mathcal {O}_S^\times \) 的截面所成的集合，与 \emph{\(\mathscr {Q}\) 的可逆商模} 所成的集合一一对应；或者，按照凝聚模 \(\mathscr {Q}\) 所伴随的射影丛 \(\mathbb {P}(\mathscr {Q})\) 的定义（参见 [\hyperref[Gro1960a]{Gro1960a}，V，§2]），与 \(\mathbb {P}(\mathscr {Q})\) 在 \(S\) 上的截面所成的集合一一对应。此描述与取逆像相容，因而得到下述定理。\begin{theorem}{4.3}\setnumber{4.3}\label{fga3.v-4-theorem-4.3}\skipsinglepar\par{}设 \(f\colon  X\to  S\) 是几何纤维整的平坦固有态射，\(S\) 局部 Noether，\(\mathscr {L}\) 是 \(X\) 上的可逆模。对每个 \(S\) 上的 \(S'\)，以 \(T(S')\) 表示 \(X\times _S S'/S'\) 上满足 \(\mathscr {L}(D)\) 在 \(S'\) 上局部同构于 \(\mathscr {L}\otimes _{\mathcal {O}_S}\mathcal {O}_{S'}\) 的相对除子 \(D\) 所成的集合（即要求 \(\mathscr {L}(D)\) 与 \(\mathscr {L}\otimes _{\mathcal {O}_S}\mathcal {O}_{S'}\) 定义 \(\operatorname {Pic}(X\times _S S'/S')\) 的同一元素）。那么，\(S\) 上存在一个在唯一同构意义下确定的凝聚模 \(\mathscr {Q}\)，使得函子 \(T\) 由射影丛 \(\mathbb {P}(\mathscr {Q})\) 表示。\end{theorem}\begin{corollary}{4.4}\setnumber{4.4}\label{fga3.v-4-corollary-4.4}\skipsinglepar\par{}若假设 \(f\) 射影，则函子性同态 \(\mathscr {D}\kern  -1pt\mathit {iv}_{X/S}\to \mathscr {P}\kern  -1pt\mathit {ic}_{X/S}\) 可由射影态射表示。\end{corollary}\par{}若 \(X\) 在 \(S\) 上有截面（分别地，局部有截面），则由 \hyperref[fga3.v-4-theorem-4.3]{定理~\ref*{fga3.v-4-theorem-4.3}} 及 \hyperref[fga3.v-2-proposition-2.1]{命题~\ref*{fga3.v-2-proposition-2.1}}，上述同态可由射影丛（分别地，局部射影丛）表示。当 \(f\) 拟射影时，利用 \(X\) 在 \(S\) 上的有限平坦局部拟截面，通过下降方法很容易归约到上述情形。\begin{remark}{4.5}\setnumber{4.5}\label{fga3.v-4-remark-4.5}\skipsinglepar\par{}在 \hyperref[fga3.v-4-theorem-4.3]{定理~\ref*{fga3.v-4-theorem-4.3}} 的条件下，\(S\) 上的模 \(\mathscr {Q}\) 一般并不局部自由；这是因为 \(\mathscr {Q}\) 的约化纤维的维数，即当 \(s\in  S\) 变化时 \(\operatorname {H}^0(X_s,\mathscr {L}_s)\) 的维数，可以 \emph{跳跃}。给定局部 Noether 预概形 \(S\) 上的凝聚模 \(\mathscr {Q}\)，容易证明：对给定的任意 \(s\in  S\)，\(\mathscr {Q}\) 在 \(s\) 处自由，当且仅当 \(\mathbb {P}(\mathscr {Q})\) 在 \(s\) 上方各点处在 \(S\) 上平坦（此时它甚至在 \(s\) 上方各点处在 \(S\) 上是光滑的）。当这一条件满足，且 \(\mathscr {Q}\) 如上由 \(\mathscr {L}\) 定义时，它还蕴含直像 \(f_*(\mathscr {L})\) 的构造 \emph{“与 \(s\) 的某个邻域上的基变换交换”}，或蕴含 \(f_*(\mathscr {L})_s\to \operatorname {H}^0(X_s,\mathscr {L}_s)\) 满射。例如，若 \(\operatorname {H}^1(X_s,\mathscr {L}_s)=0\)，就属于这一情形。只要所讨论的预概形存在，这些判据尤其适用于泛情形 \(\underline {\operatorname {Div}}_{X/S}\to \underline {\operatorname {Pic}}_{X/S}\)，并给出此态射在 \(\underline {\operatorname {Div}}_{X/S}\) 的给定点处光滑的充要条件（分别地，充分条件）。\end{remark}\chapter{主要存在定理的证明}\thispagestyle{fancy}\label{fga3.v-5}\par{}在 \hyperref[fga3.v-3-theorem-3.1]{定理~\ref*{fga3.v-3-theorem-3.1}} 的条件下，选取一个在 \(X/S\) 上甚丰沛的模 \(\mathcal {O}_X(1)\)，以 \(\xi \) 表示 \(\operatorname {Pic}(X/S)\) 中的相应元素。为简便起见，令 \(\mathscr {P}(S')=\operatorname {Pic}(X\times _S S'/S')\)，并假设 \(X/S\) 有截面。令 \(\mathscr {P}^+(S')\) 为 \(\mathscr {P}(S')\) 的如下子集：其元素是 \(X\times _S S'\) 上满足 \[     \begin {aligned}       \operatorname {R}^i f'_*(\mathscr {L}(n)) &= 0       \qquad \text {当 }i>0\text { 且任意 }n\geqslant 0     \\f'_*(\mathscr {L}(n)) &\neq 0       \qquad \text {对所有 }n\geqslant 0.     \end {aligned}   \] 的可逆模 \(\mathscr {L}\) 的类。这些条件在基变换下稳定，因而定义了 \(\mathscr {P}\) 的子函子 \(\mathscr {P}^+\)；它显然在由 \(\xi \) 给出的平移下稳定。利用 Serre 的“定理 A 和 B”（[\hyperref[GD1960]{GD1960}，III，§2]）及一般理论（[\hyperref[Gro1960a]{Gro1960a}，IV，5]），容易看出，\(\mathscr {P}\) 可表当且仅当 \(\mathscr {P}^+\) 可表；于是，\(\mathscr {P}^+\) 将由表示 \(\mathscr {P}\) 的预概形 \(\underline {\operatorname {Pic}}_{X/S}\) 中的一个开集 \(U\) 表示，而后者将是开集 \(U\setminus  n\xi \) 的递增并。\par{}为简便起见，令 \(\mathscr {D}=\mathscr {D}\kern  -1pt\mathit {iv}_{X/S}\)，并令 \(\mathscr {D}^+\) 为 \(\mathscr {P}^+\) 在典范态射 \(\mathscr {D}\to \mathscr {P}\) 下的逆像。于是有态射 \begin{eqn}{+}\setnumber{+}\label{fga3.v-5-equation-plus}\skipsinglepar\[       \mathscr {D}^+       \to  \mathscr {P}^+     \tag{+}     \]\end{eqn}；我们已经知道，\(\mathscr {D}^+\) 可由预概形 \(D=\underline {\operatorname {Div}}_{X/S}\) 中的开集 \(D^+\) 表示（更确切地说，可由处处非零的 \emph{局部自由模所伴随的射影丛} 表示）。这是因为，若 \(X\times _S S'\) 上的 \(\mathscr {L}\) 满足本节开头的条件，则 \(f'_*(\mathscr {L})\) 是非零的 \emph{局部自由} 模，其构造与基变换交换；采用 \hyperref[fga3.v-4-theorem-4.3]{定理~\ref*{fga3.v-4-theorem-4.3}} 的记号，此时 \(\mathscr {Q}\) 同构于 \(f'_*(\mathscr {L})\) 的对偶。因此，利用一般判据（[\hyperref[Gro1960a]{Gro1960a}，IV，4.7]），便可断定 \(P^+\) 可表。在 \emph{同上} 中，取 \(\mathcal {S}\) 为所有忠实平坦拟紧预概形态射所成的集合（由 \hyperref[fga3.i-b.1]{FGA 3.I, §B}，它们确为有效满态射）。正如我们刚才所见，\emph{同上} 的条件 (a)，即 \hyperref[fga3.v-5-equation-plus]{式~\ref*{fga3.v-5-equation-plus}} 可由属于 \(\mathcal {S}\) 的态射表示，已得到满足；条件 (b) 是说函子 \(\mathscr {P}^+\) 与忠实平坦拟紧下降相容，这一点是直接的。只须再证明 \emph{同上} 的条件 (c)：预概形 \(D^+\) 中由 \hyperref[fga3.v-5-equation-plus]{式~\ref*{fga3.v-5-equation-plus}} 内的 \(\mathcal {S}\)-可表态射诱导的等价关系 \(R\) 是 \(\mathcal {S}\)-有效的，即它有效且 \(D^+\to  D^+/R\) 属于 \(\mathcal {S}\)。为此首先注意，\(D^+\) 中对应于虚 Hilbert 多项式 \(Q\in \mathbb {Q}[t]\) 的开集 \({D^+}^Q\) 在 \(R\) 下稳定（因为 \(R\) 的纤维连通）；于是归约为对每个 \(Q\) 证明 \({D^+}^Q\) 上诱导的等价关系 \(R^Q\) 是 \(\mathcal {S}\)-有效的。但现在 \({D^+}^Q\) 拟射影，且等价关系 \(R^Q\) 射影且平坦。因此，\hyperref[fga3.iii-6-theorem-6.1]{FGA 3.III，定理 6.1} 的假设得到满足，它给出所需结论。\par{}在 \(X/S\) 未必有截面的一般情形中，利用下降技术很容易归约到上述情形；也可以对 \(\mathscr {P}^+\) 的定义作必要的修改，再重复上述证明。\begin{remarks}{5.1}\setnumber{5.1}\label{fga3.v-5-remarks-5.1}\skipsinglepar\par{}我们采用的方法，本质上是 Matsusaka 用来射影地构造 Picard 簇的方法。我们从 \hyperref[fga3.iii]{FGA~\ref*{fga3.iii}} 引用的、使我们能够取有效商的结果，也很容易由 Hilbert 概形的存在定理推出（例如参见 [Mumford–Tate seminar，1962]）。（在经典理论中，这些商是利用 Chow 坐标构造的。）注意，\(\underline {\operatorname {Pic}}_{X/S}\) 中开集 \(\underline {\operatorname {Pic}}_{X/S}^+\) 的构造，以及按照定义所讨论的可逆模的 \emph{除子} 的 Hilbert 多项式、将其分解为在 \(S\) 上拟射影的开集 \({\underline {\operatorname {Pic}}_{X/S}^+}^Q\) 的分解，都与基变换相容（这使我们能够使用下降技术）。\end{remarks}\begin{remarks}{5.2}\setnumber{5.2}\label{fga3.v-5-remarks-5.2}\skipsinglepar\par{}只要 \(f\colon  X\to  S\) 固有且平坦，且同态 \(k(s)\to \operatorname {H}^0(X_s,\mathcal {O}_{X_s})\)（\(s\in  S\)）都是同构，就有 \(\underline {\operatorname {Pic}}_{X/S}\) 存在，这并非不可能（后一条件还蕴含 \(\mathcal {O}_S\xrightarrow {\sim } f_*(\mathcal {O}_X)\)，且在任意基变换 \(S'\to  S\) 后仍然成立）。至少，在解析空间的范畴中，若进一步假定 \(f\) 射影，就可以用 [\hyperref[Gro1960a]{Gro1960a}，IX，3.1] 中阐明的一种微分方法来证明这一点（若我没有记错，这是 Chow 的方法）。在这种方法中，不再对除子概形的开集取一个 \emph{固有} 且平坦的等价关系的商，而是在 \(X\) 到 \(\mathbb {P}_S^r\) 的浸入概形中对射影群取商。利用 Mumford 关于对射影群 \hyperref[Mum1961]{[Mum1961]} 取商的结果，这种方法很可能可以推广到概形的情形；目前尚无成文的证明，除了 \(X\) 在 \(S\) 上有“许多局部截面”的情形，例如 \(X\) 在剩余域代数闭的完备局部环上可分的情形。原则上，这种方法适用范围更广，因为它也给出非分离 Picard 预概形的存在性，而在这种情形中，第一种方法必然失效。（从技术上说，困难在于：当 \(f\) 的几何纤维并非整的时，\hyperref[fga3.v-4-theorem-4.3]{定理~\ref*{fga3.v-4-theorem-4.3}} 中考虑的函子不再由射影丛 \(\mathbb {P}(\mathscr {Q})\) 本身表示，而由它的一个 \emph{开集} 表示；这就引出了对平坦但不固有的等价关系取商这一微妙问题。）\par{}\emph{[编者]} 正如我在下一次报告开头所指出的，这里提出的存在性猜想是错误的；但 Mumford 用他的方法证明了一个稍弱的定理。\end{remarks}\begin{remark}{5.3}\setnumber{5.3}\label{fga3.v-5-remark-5.3}\skipsinglepar\par{}注意，这里的证明既不需要预先构造曲线或曲线族的 Jacobian，也不需要 Abel 簇或 Abel 概形的理论。因此，它与沿 Weil 所勾画的道路展开的传统论述——例如 Lang 的书 \hyperref[Lan1959]{[Lan1959]} 或 Chevalley 的文章 \hyperref[Che1960]{[Che1960]}——有本质区别。即使对于代数闭域（譬如复数域）上非奇异曲线的 Jacobian，这里给出的构造也是唯一已知的、具有我们在 \cref{fga3.v-1} 中用作定义的极强性质的构造（这些性质本质上是 Chevalley 的性质，但允许“参数簇”带有幂零元）。Picard 概形的构造理应先于而非后于 Abel 簇理论，这一点先验地就很清楚：一般的 Picard 概形既不是 Abel 概形，也不能归约为 Abel 概形。代数闭域上奇异曲线的情形已经表明这一点；这里出现 Rosenlicht 的“广义 Jacobian”，它们并非 Abel 簇。此外，一旦有了一般的 Picard 概形理论，Abel 簇、更一般的 Abel 概形的理论就会大为简化。特别地，Abel 概形的对偶理论，尤其是 Cartier 型结果，此时在一定程度上成为形式的推论（例如参见 [Mumford–Tate seminar，1962]）。\end{remark}\begin{remarks}{5.4}\setnumber{5.4}\label{fga3.v-5-remarks-5.4}\skipsinglepar\par{}对于一条退化为奇异曲线的曲线的 Jacobian，Igusa 的“相容性原理”只有被理解为相对曲线概形 \(X/S\) 的 Picard 概形的存在定理时，才能得到恰当理解；这里的曲线在 \(S\) 上未必光滑。因此，当特殊纤维曲线整时（即用经典术语说，不可约且重数为 \(1\) 时），这是主要存在定理 \hyperref[fga3.v-3-theorem-3.1]{定理~\ref*{fga3.v-3-theorem-3.1}} 的一个特例。注意，目前对于可约的特殊纤维曲线（即使各分支的重数均为 \(1\)，即特殊纤维曲线在剩余域上可分），已知存在定理仍未涵盖，除非底环是剩余域代数闭的 \emph{完备} 离散赋值环；参见 \hyperref[fga3.v-5-remarks-5.2]{Remarks~\ref*{fga3.v-5-remarks-5.2}}。在概形理论中，要对亏格为 \(g\) 的曲线的模概形作 Baily–Satake“紧化”的几何构造，必然会遇到这一存在性问题。（目前，由 Igusa 的工作，只在 \(g=1\) 时知道这种紧化。）\end{remarks}\chapter{相对存在定理}\thispagestyle{fancy}\label{fga3.v-6}\par{}这里将概述一些有用的情形：某些 Picard 概形的存在性蕴含另一些 Picard 概形的存在性。这使我们能由主要存在定理 \hyperref[fga3.v-3-theorem-3.1]{定理~\ref*{fga3.v-3-theorem-3.1}} 推出各种其他存在定理。\begin{proposition}{6.1}\setnumber{6.1}\label{fga3.v-6-proposition-6.1}\skipsinglepar\par{}设 \(f\colon  X\to  S\) 是平坦射影态射，使其 Stein 分解 \(f=f''f'\) 中的态射 \(f'\colon  X\to  S'\) 为 \emph{平坦} 且几何纤维整（从而满足 \hyperref[fga3.v-3-theorem-3.1]{定理~\ref*{fga3.v-3-theorem-3.1}} 的假设），并使有限态射 \(f''\colon  S'\to \) 平坦。那么 \(\underline {\operatorname {Pic}}_{X/S}\) 存在，而且（用 \hyperref[fga3.ii-c.2]{FGA 3.II, §C.2} 中引入的记号）有典范同构 \[       \underline {\operatorname {Pic}}_{X/S}       \xrightarrow {\sim } \prod _{S'/S} \underline {\operatorname {Pic}}_{X/S'}.     \]。\begin{proof}\par{}为证明这一点，先建立函子同构 \[         \mathscr {P}\kern  -1pt\mathit {ic}_{X/S}         \xrightarrow {\sim } \prod _{S'/S} \mathscr {P}\kern  -1pt\mathit {ic}_{X/S'}       \]，再利用 \hyperref[fga3.v-3-theorem-3.1]{定理~\ref*{fga3.v-3-theorem-3.1}}，由此得到 \(\mathscr {P}\kern  -1pt\mathit {ic}_{X/S'}\) 可表。为保证 \(\prod _{S'/S}\underline {\operatorname {Pic}}_{X/S'}\) 存在，使用 \cref{fga3.v-3} 中所说明的 \(\underline {\operatorname {Pic}}_{X/S'}\) 的结构（它蕴含 \(\underline {\operatorname {Pic}}_{X/S'}\) 在 \(S\) 上的任一纤维的有限子集都包含在某个仿射开集中）。\end{proof}\end{proposition}\par{}例如，若概形 \(X\) 是 \(S\) 上满足 \hyperref[fga3.v-3-theorem-3.1]{定理~\ref*{fga3.v-3-theorem-3.1}} 的条件的概形 \(X_i\) 的 \emph{和}，则 \hyperref[fga3.v-6-proposition-6.1]{命题~\ref*{fga3.v-6-proposition-6.1}} 的陈述归结为平凡的陈述 \[     \underline {\operatorname {Pic}}_{X/S}     \xrightarrow {\sim } \prod _i \underline {\operatorname {Pic}}_{X_i/S}.   \]。\begin{corollary}{6.2}\setnumber{6.2}\label{fga3.v-6-corollary-6.2}\skipsinglepar\par{}设 \(f\colon  X\to  S\) 是几何纤维局部整的射影平坦态射（例如射影且正规的态射）。那么 \(\underline {\operatorname {Pic}}_{X/S}\) 存在。\begin{proof}\par{}这时，\(S'\) 是 \(S\) 的 étale 覆盖（只要 \(f\) 可分，即平坦且几何纤维约化，这一点就成立）；由 \(\underline {\operatorname {Pic}}_{X/S'}\) 的类似结构可见，\cref{fga3.v-3} 中为 \(\underline {\operatorname {Pic}}_{X/S}\) 陈述的结构定理仍然成立。\end{proof}\end{corollary}\par{}应用下降过程，可以得到一个相对存在定理；其适用范围取决于 \hyperref[fga3.i-a.2.c]{FGA 3.I, §A.2.c} 中提出的非平坦下降问题能否解决。这里我们仅说明一个特殊情形：\begin{proposition}{6.3}\setnumber{6.3}\label{fga3.v-6-proposition-6.3}\skipsinglepar\par{}设 \(f\colon  X\to  S\) 是固有态射，\(X_1\) 与 \(X_2\) 是 \(X\) 的子预概形，均在 \(S\) 上平坦，分别由凝聚理想 \(\mathscr {J}_1\) 与 \(\mathscr {J}_2\) 定义，使得 \(\mathscr {J}_1\cap \mathscr {J}_2=(0)\)，且 \(\mathcal {O}_X/(\mathscr {J}_1+\mathscr {J}_2)\) 在 \(S\) 上平坦（即 \(X\) 中作为 \(X_1\) 与 \(X_2\) 的上确界的子预概形就是 \(X\) 本身，而它们的下确界 \(Z\) 在 \(X\) 上平坦）。进一步假设，对所有 \(s\in  S\)，当 \(i=1,2\) 时同态 \(k(s)\to \operatorname {H}^0(X_{i_s},\mathcal {O}_{X_{i_s}})\) 均为双射。那么自然的函子同态 \[       \mathscr {P}\kern  -1pt\mathit {ic}_{X/S}       \to  \mathscr {P}\kern  -1pt\mathit {ic}_{X_1/S} \times  \mathscr {P}\kern  -1pt\mathit {ic}_{X_2/S}     \] 可由仿射态射表示。因此，若 \(\underline {\operatorname {Pic}}_{X_1/S}\) 与 \(\underline {\operatorname {Pic}}_{X_2/s}\) 存在，则 \(\underline {\operatorname {Pic}}_{X/S}\) 也存在，且典范态射 \[       \underline {\operatorname {Pic}}_{X/S}       \to  \underline {\operatorname {Pic}}_{X_1/S} \times  \underline {\operatorname {Pic}}_{X_2/S}     \] 是仿射的。\begin{proof}\par{}借助忠实平坦下降，可归约到 \(Z\) 在 \(S\) 上有截面的情形。这个截面在 \(S\) 上定义 \(X\)、\(X_1\) 及 \(X_2\) 的截面，并使我们能如 \hyperref[fga3.v-2-remark-2.5]{注~\ref*{fga3.v-2-remark-2.5}} 中所述消除所考虑结构的自同构。此时，证明在于注意到：给定 \(X\) 上的“刚性化”可逆模 \(\mathscr {L}\)，等价于给定三元组 \((\mathscr {L}_1,\mathscr {L}_2,u)\)，其中 \(\mathscr {L}_i\) 是 \(X_i\) 上的“刚性化”模，而 \(u\) 是从 \(\mathscr {L}_1|Z\) 到 \(\mathscr {L}_2|Z\) 的、与刚性化相容的同构。只须再验证，当 \(\mathscr {L}_1\) 与 \(\mathscr {L}_2\) 固定时，给定 \(u\) 等价于给定 \(S\) 上某个适当的、在 \(S\) 上仿射的概形的一个截面；这一点很容易。\end{proof}\end{proposition}\par{}由 \hyperref[fga3.v-6-proposition-6.3]{命题~\ref*{fga3.v-6-proposition-6.3}} 容易得出：\begin{corollary}{6.4}\setnumber{6.4}\label{fga3.v-6-corollary-6.4}\skipsinglepar\par{}设 \(X\) 是域 \(k\) 上的固有可分概形，\(X_i\) 是 \(X\) 的不可约分支。若各 \(\underline {\operatorname {Pic}}_{X_i/k}\) 存在，则 \(\underline {\operatorname {Pic}}_{X/k}\) 也存在，且典范态射 \[       \underline {\operatorname {Pic}}_{X/k}       \to  \prod _i \underline {\operatorname {Pic}}_{X_i/k}     \] 是仿射的。\end{corollary}\par{}与 \hyperref[fga3.v-6-corollary-6.2]{推论~\ref*{fga3.v-6-corollary-6.2}} 结合，这表明，例如，只要 \(X\) 是域 \(k\) 上的射影可分概形，\(\underline {\operatorname {Pic}}_{X/k}\) 就存在。即使 \(X\) 在 \(k\) 上不再可分，利用 Oort 的论证 \hyperref[Oor1962]{[Oor1962]}，同样也有归约结果。这一方法也适用于任意底上的概形（例如，在下一次报告中证明 \hyperref[fga3.v-3-remark-3.3]{注~\ref*{fga3.v-3-remark-3.3}} 所陈述的有限性结果时，这种情形就很有用）。为避免过于技术性的陈述，我们只考虑底为域的情形：\begin{proposition}{6.5}\setnumber{6.5}\label{fga3.v-6-proposition-6.5}\skipsinglepar\par{}设 \(X\) 是域 \(k\) 上的固有概形，\(X_0\) 是一个具有相同底层集合的子概形（因而由 \(X\) 上的幂零理想定义）。那么函子性态射 \(\mathscr {P}\kern  -1pt\mathit {ic}_{X/k}\to \mathscr {P}\kern  -1pt\mathit {ic}_{X_0/k}\) 可由仿射态射表示。特别地，若 \(\underline {\operatorname {Pic}}_{X_0/k}\) 存在，则 \(\underline {\operatorname {Pic}}_{X/k}\) 也存在，且态射 \(\underline {\operatorname {Pic}}_{X/k}\to \underline {\operatorname {Pic}}_{X_0/k}\) 是仿射的。\end{proposition}\par{}将它与 \hyperref[fga3.v-6-corollary-6.4]{推论~\ref*{fga3.v-6-corollary-6.4}} 结合，容易得出：\begin{corollary}{6.6}\setnumber{6.6}\label{fga3.v-6-corollary-6.6}\skipsinglepar\par{}设 \(X\) 是域 \(k\) 上的射影预概形。那么 \(\underline {\operatorname {Pic}}_{X/k}\) 存在。\end{corollary}\begin{remark}{6.6}\setnumber{6.6}\label{fga3.v-6-remark-6.6}\skipsinglepar\par{}对于域 \(k\) 上的每个 \emph{固有} 概形 \(X\)，\(\underline {\operatorname {Pic}}_{X/k}\) 应当存在，这是极为可信的。由上述结果，这个问题可以归约到 \(k\) 代数闭、\(X\) 整的情形。此时知道，存在一个在 \(k\) 上 \emph{射影} 的整概形 \(X'\) 及支配态射 \(g\colon  X'\to  X\)（Chow 引理）。因此，只须证明相应的函子性态射 \(\mathscr {P}\kern  -1pt\mathit {ic}_{X/k}\to \mathscr {P}\kern  -1pt\mathit {ic}_{X'/k}\) 可表（很可能还可由仿射态射表示），因为已经知道 \(\mathscr {P}\kern  -1pt\mathit {ic}_{X'/k}\) 可表。这引出了目前尚无法回答的非平坦下降问题。注意，若只考虑函子 \(\mathscr {P}\kern  -1pt\mathit {ic}_{X/k}\) 在约化预概形上的限制（其中 \(X\) 是代数闭域 \(k\) 上的固有整概形），则确实得到一个可表函子：在 \(X\) 正规的情形中，Chevalley 在 \hyperref[Che1960]{[Che1960]} 中证明了这一点；在一般情形中，Seshadri 在 \hyperref[Ses1962]{[Ses1962]} 中用下降方法证明了这一点。不过，按我们的记号，这些作者所构造的概形并非 \(\underline {\operatorname {Pic}}_{X/k}\)，而是 \((\underline {\operatorname {Pic}}_{X/k})_\mathrm {red}\)，即与 \(\underline {\operatorname {Pic}}_{X/k}\) 对应的约化概形。\par{}\emph{[编者]} 正如我们在 \hyperref[fga3.vi]{下一次报告} 开头所指出的，Murre 最近已经对这里提出的问题给出了肯定回答。\end{remark}\begin{remark}{6.7}\setnumber{6.7}\label{fga3.v-6-remark-6.7}\skipsinglepar\par{}更一般地，设 \(f\colon  X'\to  X\) 是 \(k\) 上固有预概形之间的满射态射。考虑非平坦下降，就会猜想 \(\underline {\operatorname {Pic}}_{X/k}\to \underline {\operatorname {Pic}}_{X'/k}\) 是 \emph{仿射} 态射；特别地，这将蕴含（除去单位元的连通分支后）相应的 Néron–Severi 群同态模去挠元后为单射。当 \(X\) 和 \(X'\) 非奇异时，利用相交理论可验证这一点。在其他情形中，答案似乎尚未知道。\par{}\emph{[编者]} 这里提出的问题已得到肯定回答（参见下一次报告的评论的最后一段）。\end{remark}\begin{remark}{6.8}\setnumber{6.8}\label{fga3.v-6-remark-6.8}\skipsinglepar\par{}考虑带有幂零元的代数概形的 Picard 概形，对于许多问题是有用的，甚至不可或缺；这或许出乎意料。例如，若 \(X\) 是 \(k\) 上的光滑射影概形，\(Y\) 是一个超平面截面，就可以考虑 \(Y\) 的各阶“无穷小邻域”\(X_n\) 及其 Picard 概形 \(\underline {\operatorname {Pic}}_{X_n/k}\)。当 \(X\) 不可约且维数为 \(\geqslant 4\)（分别地，\(\geqslant 3\)）时，典范态射 \[       \underline {\operatorname {Pic}}_{X/k}       \to \underline {\operatorname {Pic}}_{X_n/k}       \qquad \text {对充分大的 }n     \] 是同构（分别地，在 Néron–Severi 群的挠子群的逆像之间诱导同构）；下一次报告对 Picard 概形作定性研究时将用到这一结果。同样，考虑某些带有幂零元的曲线的 Picard 概形，并利用形式几何的基本定理 [Gro1958a]，就能在等特征情形中证明 Mumford 的一个猜想：对于每个维数为 \(2\) 的完备正规 Noether 局部环 \(A\)，\(A\) 的除子类群可以看成剩余域 \(k\) 上的一个代数群 \(G\) 的 \(k\)-有理点所成的集合（一旦选定 \(A\) 中的代表域，\(G\) 就典范地确定）。当 \(A\) 的维数任意时，似乎应当存在一个 \(k\) 上的代数 pro-群，起到上述 \(G\) 的作用；当可以对 \(\operatorname {Spec}(A)\) 作“解奇点”时，它可构造成 \(k\) 上适当的（带有幂零元的）射影概形的 Picard 概形的投射极限。\end{remark}
\section*{原报告参考文献}
\begin{enumerate}
\item[{[1]}] Chevalley (Claude). \emph{Sur la théorie de Picard}. Amer. J. of Math., t.~82, 1960, pp.~435--490.
\item[{[2]}] Grothendieck (A.) et Dieudonné (J.). \emph{Éléments de géométrie algébrique. Chapitre I et suivants}. Paris, Presses universitaires de France, 1960--1961 (Institut des Hautes Études Scientifiques, Publications mathématiques 4, 8, \ldots).
\item[{[3]}] Grothendieck (Alexander). \emph{Géométrie formelle et géométrie algébrique}. Séminaire Bourbaki, t.~11, 1958/59, n\textsuperscript{o}~182, 28~页。
\item[{[4]}] Grothendieck (Alexander). \emph{Techniques de descente et théorèmes d'existence en géométrie algébrique, I--IV}. Séminaire Bourbaki, t.~12, 1959/60, n\textsuperscript{o}~190, 29~页；n\textsuperscript{o}~195, 22~页；t.~13, 1960/61, n\textsuperscript{o}~212, 20~页；n\textsuperscript{o}~221, 28~页。
\item[{[5]}] Grothendieck (Alexander). \emph{Techniques de construction en géométrie analytique, I--X}. Séminaire Cartan, t.~13, 1960/61: \emph{Déformation des structures analytiques complexes}, n\textsuperscript{o}~7--17.
\item[{[6]}] Grothendieck (Alexander). \emph{Séminaire de géométrie algébrique}，由 Lichtenbaum 撰写，Harvard University, 1961（待出版）。
\item[{[7]}] Lang (Serge). \emph{Abelian varieties}. New York, Interscience Publishers, 1959 (Interscience Tracts in Pure and Applied Mathematics, 7).
\item[{[8]}] Mumford (D.). \emph{An elementary theorem in geometric invariant theory}. Bull. Amer. Math. Soc., t.~67, 1961, pp.~483--486.
\item[{[9]}] Mumford (D.). \emph{The topology of normal singularities of an algebraic surface and a criterion for simplicity}. Paris, Presses universitaires de France, 1961 (Institut des Hautes Études Scientifiques, Publications mathématiques, 9).
\item[{[10]}] Mumford (D.) and Tate (J.). \emph{Séminaire de géométrie algébrique}. Harvard University, Spring Term 1962（待出版）。
\item[{[11]}] Oort (F.). \emph{Sur le schéma de Picard}. Bull. Soc. Math. France, t.~90, 1962（待出版）。
\item[{[12]}] Seshadri (C. S.). 学位论文，待发表于 \emph{Annali di Matematica}。
\end{enumerate}
